\sscp{Hoti}{10-03-05}
Hoti is a (probably extinct) language of northeast Seram spoken
in the village of the same name.
It is known from a word list in \citet{Ludeking1868}.
Based on this data \citet{Stresemann1927} included Hoti in his Sub-Seran subgroup.
\citet[117]{Collins1982c} writes “Hoti
and Liambata, as Stresemann’s record of them can be interpreted,
are Seti languages [{\ldots}]”.
However, the data which is available is very different from \tsc{Seti}
(\srf{10-03-04}) and Hoti is clearly a separate language based on different sound
correspondences and the lexicon.
Hoti has a large number of unconditioned splits which makes classifying it difficult.

Our discussion is based primarily on the data as presented by
\citeauthor{Stresemann1927}’s \citeyear{Stresemann1927} interpretation of \citet{Ludeking1868}.
Data taken directly from \citet{Ludeking1868} are indicated.
We interpret ‹ph› as a bilabial fricative [ɸ].
PAMS reconstructions in this section follow Stresemann’s transcription.

We place Hoti in \tsc{Ambon-Seram},
primarily based on the merger of *n/*ŋ,
combined with the lack of any positive evidence for grouping
it with any branch of \tsc{Seram-Tanimbar-Bomberai}.
We further identify it as a primary branch of \tsc{Ambon-Seram}
due to the fact that *R does not fully merge with *l.

Hoti appears to merge *n/*ñ/*ŋ,
with a subsequent split to \it{n, l, r}.
The outcome of the split is partially conditioned by the quality
of the following vowel.
Before a high vowel *n/*ñ/*ŋ > \it{r},
with only one instance of *ŋ > \it{l} in this environment.
Examples are given in \qf{10-71}.

\noindent{\wg\st{6pt}\tblr{>{\hspace{-\tabcolsep}}l<{\hspace{-\tabcolsep}}
>{\hspace{-\tabcolsep}\enspace}
lllll}
\lsptoprule
%\tbx{10-71}	&	*n/*ŋ	&	env.	&	PMP/PAMS	&	Hoti	&	local gl.	\\	\midrule
%\endfirsthead
\tbx{10-71}	&	**n	&	env.	&	PMP/PAMS	&	Hoti	&	local gl.	\\	\midrule
%\endhead
	&	{\hr\hr}l	&	{\#}{\gap}V \tsc{[+h]}	&	*\tbr{ŋ}is(ŋ)i(s)	&	\tbr{l}isi-	&	tooth	\\
	&	{\hr\hr}r	&	{\#}{\gap}V \tsc{[+h]}	&	*\tbr{ñ}iuR	&	\tbr{r}uh-a	&	coconut	\\
	&	{\hr\hr}r	&	{\#}{\gap}V \tsc{[+h]}	&	**\tbr{n}iabət	&	\tbr{r}iapat-a	&	flying fox	\\
	&	{\hr\hr}r	&	{\#}{\gap}V \tsc{[+h]}	&	**\tbr{n}ipil	&	\tbr{r}in-a	&	bee	\\
	&	{\hr\hr}r	&	{\#}{\gap}V \tsc{[+h]}	&	*\tbr{ŋ}ijuŋ	&	\tbr{r}iri-	&	nose	\\
	&	{\hr\hr}r	&	V{\gap}V \tsc{[+h]}	&	*ma-h\<i\tbr{n}\>ipi >	&	mi-\tbr{r}ii	&	dream	\\	\hhline{~}
	&						&											&	**ma\tbr{n}ipi	&		&		\\
	&	{\hr\hr}r	&	V{\gap}V \tsc{[+h]}	&	*ma\tbr{n}uk	&	ma\tbr{r}iʔ-a	&	bird	\\
\lspbottomrule
\end{tabular}\mbox{}\\}

Before non-high vowels the reflexes are more diverse.
*n/*ñ/*ŋ > \it{r} is only attested after high vowels,
(three examples), but *n/*ŋ > \it{n,
l} also occur in this environment (two examples each).
Word initially before a non-high vowel,
or word medially between non-high vowels,
*n/*ŋ > \it{n, l.} Examples are given in \qf{10-72}.

{\wg\st{4pt}\begin{xltabular}[l]{\linewidth}
{>{\hspace{-\tabcolsep}}l<{\hspace{-\tabcolsep}}
>{\hspace{-\tabcolsep}\enspace}llll
>{\raggedright\arraybackslash}X}
\lsptoprule
\tbx{10-72}		&	**n	&	env.	&	PMP/PAMS	&	Hoti	&	local gl.	\\	\midrule	\hhline{~}
\endfirsthead
\lsptoprule
\qf{10-72}	{\cwh}&	**n	{\cwh}&	env.	{\cwh}&	PMP/PAMS	{\cwh}&	Hoti	{\cwh}&	local gl.	{\cwh}\\	\midrule	\hhline{~}
\endhead
	&	{\hr\hr}n	&	{\gap}V\tsc{[-h]}	&	*\tbr{n}anaq	&	\tbr{n}ana	&	pus	\\
	&	{\hr\hr}n	&	{\gap}V\tsc{[-h]}	&	*na\tbr{n}aq	&	na\tbr{n}a	&	pus	\\
	&	{\hr\hr}n	&	{\gap}V\tsc{[-h]}	&	*\tbr{ñ}awa	&	\tbr{n}awa	&	breathe	\\
	&	{\hr\hr}n	&	{\gap}C	&	*ku\tbr{n}ij	&	i\tbr{n}n-a	&	turmeric	\\
	&	{\hr\hr}n	&	{\gap}V\tsc{[-h]}	&	*ma-b‹i\tbr{n}›ahi >&	ma-hi\tbr{n}a(-a)	&	woman	\\	\hhline{~}
	&						&											&	**mabi\tbr{n}a	&		&		\\
	&	{\hr\hr}l	&	{\gap}V\tsc{[-h]}	&	*\tbr{n}aRah/*\tbr{n}ara	&	\tbr{l}aa	&	\it{Pterocarpus indicus} (k.o. tree)	\\
	&	{\hr\hr}l	&	{\gap}V\tsc{[-h]}	&	*tiRəm >	&	tiha\tbr{l}-a	&	oyster	\\	\hhline{~}
	&						&											&	**tiRə\tbr{n}	&		&		\\
	&	{\hr\hr}l	&	{\gap}V\tsc{[-h]}	&	*a\tbr{n}ak	&	a\tbr{l}a-	&	child	\\
	&	{\hr\hr}l	&	{\gap}V\tsc{[-h]}	&	*qaluhipa\tbr{n} >	&	uria\tbr{l}-a	&	centipede	\\ \hhline{~}
	&						&											&	**qalipa\tbr{n}	&		&		\\
	&	{\hr\hr}l	&	{\gap}V\tsc{[-h]}	&	*pə\tbr{ñ}u	&	he\tbr{l}-a	&	turtle	\\
	&	{\hr\hr}l	&	{\gap}V\tsc{[-h]}	&	**bərə\tbr{n}	&	bohi\tbr{l}-a	&	nipa palm	\\
	&	{\hr\hr}l	&	{\gap}V\tsc{[-h]}	&	Mly. \it{duria\tbr{n}}	&	ru\tbr{l}-a	&	durian	\\
	&	{\hr\hr}r	&	{\gap}V\tsc{[-h]}	&	**maleu\tbr{n}	&	mi-rau\tbr{r}-a	&	bush turkey	\\
	&	{\hr\hr}r	&	{\gap}V\tsc{[-h]}	&	*duyu\tbr{ŋ}	&	rui\tbr{r}-a	&	dugong	\\
	&	{\hr\hr}r	&	{\gap}V\tsc{[-h]}	&	*tali\tbr{ŋ}a	&	ti\tbr{r}a	&	ear	\\
	&	{\hr\hr}r	&	{\gap}V\tsc{[-h]}	&	*mi\tbr{ñ}ak	&	mi\tbr{r}a	&	fat (verb/adjective)	\\
\lspbottomrule
\end{xltabular}\addtocounter{table}{-1}}

In word-final position *n/*ñ/*ŋ > \it{n} in all seven instances.
Examples are given in \qf{10-73}.
Note that while there are some examples of word-final *n/*ñ/*ŋ >
\it{l} in \qf{10-72} above,
all occur before a suffix,
unlike those in \qf{10-73}.

\noindent{\wg\st{6pt}\tblr{>{\hspace{-\tabcolsep}}l<{\hspace{-\tabcolsep}}
>{\hspace{-\tabcolsep}\enspace}
lllll}
\lsptoprule
%\tbx{10-73}	&	*n/*ŋ	&	env.	&	PMP/PAMS	&	Hoti	&	local gl.	\\	\midrule
%\endfirsthead
\tbx{10-73}	&	**n	&	env.	&	PMP/PAMS	&	Hoti	&	local gl.	\\	\midrule
%\endhead
	&	{\hr\hr}n	&	{\gap\#}	&	*hika\tbr{n}	&	iʔa\tbr{n}	&	fish	\\
	&	{\hr\hr}n	&	{\gap\#}	&	**avu\tbr{n}	&	ʤafi\tbr{n}	&	fruit dove	\\
	&	{\hr\hr}n	&	{\gap\#}	&	*ka-wana\tbr{n}	&	a-wala\tbr{n}	&	right	\\
	&	{\hr\hr}n	&	{\gap\#}	&	**be\tbr{n}e	&	be{\tl}be\tbr{n}	&	jam, plug, congest	\\
	&	{\hr\hr}n	&	{\gap\#}	&	*kaə\tbr{n}	&	a\tbr{n}	&	eat	\\
	&	{\hr\hr}n	&	{\gap\#}	&	*ta\tbr{ŋ}is	&	ʤa\tbr{n}	&	cry	\\
	&	{\hr\hr}n	&	{\gap\#}	&	*i\tbr{n}um	&	i\tbr{n}	&	drink	\\
\lspbottomrule
\end{tabular}\mbox{}\\}

The consonants *d
and *l also merge.
Before high vowels the usual reflex is \it{r},
though there is also one instance each of *d > \it{l}
and *l = \it{l} before high vowels.
Examples are given in \qf{10-74}.
\vspace{2pt}

\noindent{\wg\st{6pt}\begin{tabularx}
{\linewidth}{>{\hspace{-\tabcolsep}}l<{\hspace{-\tabcolsep}}
>{\hspace{-\tabcolsep}\enspace}llll
>{\raggedright\arraybackslash}X}
\lsptoprule
%\tbx{10-74}	&	*d/*l	&	env.	&	PMP/PAMS	&	Hoti	&	local gl.	\\	\midrule
%\endfirsthead
\tbx{10-74}	&	*d/*l	&	env.	&	PMP/PAMS	&	Hoti	&	local gl.	\\	\midrule
%\endhead
	&	{\hr}r	&	{\gap}V\tsc{[+h]}	&	*\tbr{d}uyuŋ	&	\tbr{r}uir-a	&	dugong	\\
	&	{\hr}r	&	{\gap}V\tsc{[+h]}	&	Mly. \tbr{d}urian	&	\tbr{r}ul-a	&	durian	\\
	&	{\hr}r	&	{\gap}V\tsc{[+h]}	&	*\tbr{l}iqə	&	\tbr{r}ia	&	voice	\\
	&	{\hr}r	&	{\gap}V\tsc{[+h]}	&	**\tbr{l}usi	&	\tbr{r}is-a	&	hawk	\\
	&	{\hr}r	&	{\gap}V\tsc{[+h]}	&	**ta\tbr{d}igi	&	a-ka\tbr{r}ihi	&	fingernail	\\
	%&	{\hr}r	&	{\gap}V\tsc{[+h]}	&	*qaluhipan > 	&	u\tbr{r}ial-a	&	centipede	\\ \hhline{~}
	%&					&										&	**qa\tbr{l}ipan	&		&		\\
	&	{\hr}r	&	{\gap}V\tsc{[+h]}	&	**qa\tbr{l}ipan 	&	u\tbr{r}ial-a	&	centipede	\\
	&	{\hr}r	&	{\gap}V\tsc{[+h]}	&	*qu\tbr{l}u	&	u\tbr{r}i-	&	head	\\
	&	{\hr}r	&	{\gap}V\tsc{[+h]}	&	**bu\tbr{l}ut	&	bu\tbr{r}it	&	glue	\\
	&	{\hr}r	&	{\gap}V\tsc{[+h]}	&	*bu\tbr{l}an	&	ɸu\tbr{r}in	&	month	\citep{Ludeking1868}\\
	&	{\hr}l	&	{\gap}V\tsc{[+h]}	&	*\tbr{d}ilaq	&	\tbr{l}ila-	&	tongue (progressive ass.?)	\\
	&	{\hr}l	&	{\gap}V\tsc{[+h]}	&	*\tbr{l}asə	&	\tbr{l}isi-n	&	testicles	\\
\lspbottomrule
\end{tabularx}\mbox{}\\}

Before non-high vowels the usual reflex is *d/*l > \it{l},
but there are also three instances of *d/*l > \it{n} (one in
absolute word-final position,
two morpheme finally before nominal \it{-a})
and two instances of *l > \it{r} in this environment.
The quality of the previous vowel does not seem to be a conditioning
factor for *d/*l, though there are only three examples of *d/*l after a high vowel
and before a non-high vowel (/V \tsc{[+h]}{\gap}V\tsc{[-h]}).
Examples are given in \qf{10-75}.

\noindent{\wg\st{6pt}\tblr{>{\hspace{-\tabcolsep}}l<{\hspace{-\tabcolsep}}
>{\hspace{-\tabcolsep}\enspace}
lllll}
\lsptoprule
%\tbx{10-75}	&	*d/*l	&	environment	&	PMP/PAMS	&	Hoti	&	local gl.	\\	\midrule
%\endfirsthead
\tbx{10-75}	&	*d/*l	&	environment	&	PMP/PAMS	&	Hoti	&	local gl.	\\	\midrule
%\endhead
	&	{\hr}l	&	{\gap}V\tsc{[-h]}	&	*\tbr{d}aləm	&	\tbr{l}ari-	&	inside	\\
	&	{\hr}l	&	{\gap}V\tsc{[-h]}	&	*\tbr{d}aRaq	&	\tbr{l}aha	&	blood	\\
	&	{\hr}l	&	{\gap}V\tsc{[-h]}	&	*\tbr{d}ədap	&	\tbr{l}ol-a	&	Erythrina (tree)	\\
	&	{\hr}l	&	{\gap}V\tsc{[-h]}	&	*kə\tbr{d}əŋ	&	o\tbr{l}o	&	stand	\\
	&	{\hr}l	&	{\gap}V\tsc{[-h]}	&	**ma\tbr{nt}əR >	&	ma\tbr{l}a	&	cuscus	\\
	&					&										&	**ma\tbr{d}əR	&		&		\\
	&	{\hr}l	&	{\gap}V\tsc{[-h]}	&	*də\tbr{d}ap	&	lo\tbr{l}-a/lo\tbr{l}a	&	Erythrina (tree)	\\
	&	{\hr}l	&	{\gap}V\tsc{[-h]}	&	**tave\tbr{d}a	&	taɸi\tbr{l}a	&	jackfruit	\\
	&	{\hr}l	&	{\gap}V\tsc{[-h]}	&	*\tbr{l}awaq	&	\tbr{l}awa-lawa	&	spider	\\
	&	{\hr}l	&	{\gap}V\tsc{[-h]}	&	**\tbr{l}osa	&	\tbr{l}osa	&	calcium	\\
	&	{\hr}l	&	{\gap}V\tsc{[-h]}	&	*di\tbr{l}aq	&	li\tbr{l}a-	&	tongue	\\
	&	{\hr}l	&	{\gap}V\tsc{[-h]}	&	**ke\tbr{l}e	&	e\tbr{l}e	&	flank	\\
	&	{\hr}l	&	{\gap}V\tsc{[-h]}	&	**təbə\tbr{l}ə	&	tiba\tbr{l}a	&	snake	\\
	&	{\hr}l	&	{\gap}V\tsc{[-h]}	&	**və\tbr{l}ə	&	ɸe\tbr{l}aʔ-a	&	oil	\\
	&	{\hr}l	&	{\gap}V\tsc{[-h]}	&	*ka\tbr{l}aw	&	ka\tbr{l}a	&	hornbill	\\
	&	{\hr}n	&	{\gap\#}-a				&	**nipi\tbr{l}	&	ri\tbr{n}-a	&	bee	\\
	&	{\hr}n	&	{\gap\#}-a				&	**tabə\tbr{l}ə	&	tibo\tbr{n}-a	&	bamboo	\\
	&	{\hr}n	&	{\gap\#}					&	**we\tbr{d}i	&	we\tbr{n}	&	sword grass	\\
	&	{\hr}r	&	{\gap}V\tsc{[-h]}	&	**ma\tbr{l}eun	&	mi-\tbr{r}aur-a	&	bush turkey (?assimilation)	\\
	&	{\hr}r	&	{\gap}V\tsc{[-h]}	&	**ma-sə\tbr{l}ə	&	miso\tbr{r}a	&	sick	\\
\lspbottomrule
\end{tabular}\mbox{}\\}

\newpage
There are four instances of *j in available data.
Two have *j > \it{r},
one has *j > \it{n} word finally,
and one has *j > {\0} word finally.
These examples are given in \qf{10-76}.
There is one instance of *z in the available data: *qazay > \it{ala-} ‘chin’.

\noindent{\wg\st{6pt}\tblr{>{\hspace{-\tabcolsep}}l<{\hspace{-\tabcolsep}}
>{\hspace{-\tabcolsep}\enspace}
lllll}
\lsptoprule
%\tbx{10-76}	&	*j	&	environment	&	PMP/PAMS	&	Hoti	&	gloss	\\	\midrule
%\endfirsthead
\tbx{10-76}	&	*j	&	environment	&	PMP	&	Hoti	&	gloss	\\	\midrule
%\endhead
	&	{\hr}r	&	{\gap}V \tsc{[+h]}	&	*ŋi\tbr{j}uŋ	&	ri\tbr{r}i-	&	nose	\\
	&	{\hr}r	&	{\gap}V \tsc{[-h]}	&	*pa\tbr{j}ay	&	ha\tbr{r}a	&	rice	\\
	&	{\hr}n	&	{\gap}{\#}	&	*kuni\tbr{j}	&	in\tbr{n}-	&	turmeric	\\
	&	{\hr}{\0}	&	{\gap}{\#}	&	*pusə\tbr{j}	&	use-	&	navel	\\
\lspbottomrule
\end{tabular}\mbox{}\\}

PMP *R
and PAMS *r show a split in Hoti with *R > \it{h,} {\0}\it{, l}.
There are no instances of *R > \it{r} in the available data,
but this may be due to the lack of sufficient examples before a high vowel.
Examples are given in \qf{10-77}.
Note that nearby \tsc{Seti} also has *R > \it{h} as a minority
pattern \srf{10-03-04}).

\noindent{\wg\st{4pt}\tblr{>{\hspace{-\tabcolsep}}l<{\hspace{-\tabcolsep}}
>{\hspace{-\tabcolsep}\enspace}
lllll}
\lsptoprule
%\tbx{10-77}	&	*R	&	PMP/PAMS	&	Hoti	&	gloss	\\	\midrule
%\endfirsthead
\tbx{10-77}	&	*R	&	PMP/PAMS	&	Hoti	&	gloss	\\	\midrule
%\endhead
	&	{\hr}h	&	*da\tbr{R}aq	&	la\tbr{h}a	&	blood	\\
	&	{\hr}h	&	*ti\tbr{R}ən	&	ti\tbr{h}al-a	&	oyster	\\
	&	{\hr}h	&	**bə\tbr{r}ən	&	bo\tbr{h}il-a	&	Nipa palm	\\
	&	{\hr}h	&	*ñiu\tbr{R}	&	ru\tbr{h}-a	&	coconut	\\
	&	{\hr}h	&	*pa\tbr{R}ih	&	a\tbr{h}-a	&	ray	\\
	&	{\hr}h	&	**soku\tbr{r}	&	su\tbr{h}-a	&	galangal	\\
	&	{\hr}h	&	**wa\tbr{r}u	&	i-wa\tbr{h}u-a	&	k.o. fish	\\
	&	{\hr}l	&	*wahiyə\tbr{R}	&	wa\tbr{l}-a	&	water	\\
	&	{\hr}l	&	*qaba\tbr{R}a	&	ɸa\tbr{l}a-	&	front limbs	\\
	&	{\hr}l	&	*ba\tbr{R}u	&	ɸa\tbr{l}-a	&	hibiscus	\\
	&	{\hr}l	&	**kasawa\tbr{r}i	&	asaa\tbr{l}-a	&	cassowary	\\
	&	{\hr}{\0}	&	Mly. \it{du\tbr{r}ian} > **ruial	&	rul-a	&	durian	\\
	&	{\hr}{\0}	&	*na\tbr{R}a/*na\tbr{r}a	&	laa	&	\it{Pterocarpus indicus} (k.o. tree)	\\
	&	{\hr}{\0}/ʔ	&	*u\tbr{R}at	&	uat-a / uʔat-a 	&	vein, cord	\\
\lspbottomrule
\end{tabular}\mbox{}\\}

The different reflexes of Hoti *d,
*l, *R, *j, *n, *ñ,
and *ŋ are summarised in \qf{10-78}.
Minority reflexes with only a small number of examples are given in brackets.
(There are not enough examples of *j
and *z to identify any definitive patterns).
While we can posit a complete merger of *d/*l (which further
possibly merge with *n/*ŋ),
we can only posit a partial merger of *R with these segments.
Thus, if Hoti is indeed an \tsc{Ambon-Seram} language,
it does not attest the merger of *R/*l seen in all other branches except Boano (\srf{10-03-01}).
Hoti may thus be another primary branch of \tsc{Ambon-Seram}.
We note also that the (partial) conditioning of *d/*l/*n/*ŋ >
\it{l, r} according to the height of adjacent vowels is a pattern seen
only in \tsc{Ambon-Seram} languages
and is not found in \tsc{Seram-Tanimbar-Bomberai}.

\noindent{\wg\st{6pt}\tblr{>{\hspace{-\tabcolsep}}l<{\hspace{-\tabcolsep}}
>{\hspace{-\tabcolsep}\enspace}
lrlll}
\tbx{10-78}	&	a.	&	*n/*ñ/*ŋ	&	>	&	\it{r, (l)}		&	/{\gap}V\tsc{[+high]}	\\
						&			&						&	>	&	\it{l, n, r}	&	/{\gap}V\tsc{[-high]}	\\
						&			&						&	>	&	\it{n}				&	/{\gap}{\#}	\\
						&	b.	&	*d/*l			&	>	&	\it{r, (l)}		&	/{\gap}V\tsc{[+high]}	\\
						&			&						&	>	&	\it{l, (r)}		&	/{\gap}V\tsc{[-high]}	\\
						&			&						&	>	&	\it{n}				&	/{\gap}(-a){\#}	\\
						&	c.	&	*j				&	>	&	\it{r}				&		\\
						&			&						&	>	&	\it{n,} {\0}	&	/{\gap}{\#}	\\
						&	d.	&	*R				&	>	&	\it{h,} {\0}, \it{l}	&		\\
\end{tabular}\mbox{}\\}

Consistent with other languages of eastern Seram,
the usual reflex of penultimate *ə in Hoti is \it{o}.
This occurs in \tsc{Patakai-Manusela} (\srf{10-03-01})
and \tsc{Seti} (\srf{10-03-04}) to the west,
as well as the \tsc{East Seram} node of \tsc{Seram-Tanimbar-Bomberai}
(\srf{11-02}) to the south and east.
There are also two instances of penultimate *ə > \it{e},
and one instance of *ə > \it{a}.
Examples are given in \qf{10-79}.

\noindent{\wg\st{6pt}\tblr{>{\hspace{-\tabcolsep}}l<{\hspace{-\tabcolsep}}
>{\hspace{-\tabcolsep}\enspace}
lllll}
\lsptoprule
%\tbx{10-79}	&	*ə	&	PMP/PAMS	&	Hoti	&	gloss	\\	\midrule
%\endfirsthead
\tbx{10-79}	&	*ə	&	PMP/PAMS	&	Hoti	&	gloss	\\	\midrule
%\endhead
	&	{\hr}o	&	**b\tbr{ə}rən	&	b\tbr{o}hil-a	&	\it{Nipa} palm	\\
	&	{\hr}o	&	*k\tbr{ə}dəŋ	&	\tbr{o}lo	&	stand	\\
	&	{\hr}o	&	*d\tbr{ə}dap	&	l\tbr{o}l-a / lola	&	\it{Erythrina} tree	\\
	&	{\hr}o	&	**tab\tbr{ə}lə	&	tib\tbr{o}n-a	&	bamboo	\\
	&	{\hr}o	&	*k\tbr{ə}təp	&	\tbr{o}ko	&	bite	\\
	&	{\hr}o	&	*q\tbr{ə}tut	&	\tbr{o}kit	&	fart	\\
	&	{\hr}o	&	**ma-s\tbr{ə}lə	&	mis\tbr{o}ra	&	sick	\\
	&	{\hr}o	&	*t\tbr{ə}buh	&	t\tbr{o}ɸ-a	&	sugarcane	\\
	&	{\hr}o	&	*b\tbr{ə}qbəq	&	ɸ\tbr{o}ɸo-n	&	mouth \citep{Ludeking1868}	\\
	&	{\hr}o	&	*t\tbr{ə}ktək	&	t\tbr{o}to	&	cut off	\\
	&	{\hr}e	&	**v\tbr{ə}lə	&	ɸ\tbr{e}laʔ-a	&	oil	\\
	&	{\hr}e	&	*p\tbr{ə}ñu	&	h\tbr{e}l-a	&	turtle	\\
	&	{\hr}a	&	**təb\tbr{ə}lə	&	tib\tbr{a}la	&	snake	\\
\lspbottomrule
\end{tabular}\mbox{}\\}

In the final syllable the most common reflex of *ə is \it{a} (six examples).
There are also four examples of final syllable *ə > \it{o},
and *ə > \it{i}.
The larger number of  examples of *ə > \it{a} in final syllables
may point to a connection with \tsc{Seram-Tanimbar-Bomberai}
(\crf{ch:STB}) in which *ə > \it{a} in final syllables is one
of the defining changes.
Note, however, that in the two words with cognates in Bobot
and Masiwang -- the members \tsc{Seram-Tanimbar-Bomberai} closest
to Hoti -- the reflexes of schwa in final syllables do not agree.
Examples of schwa in final syllables are given in \qf{10-80}.
\\

\noindent{\wg\st{6pt}\begin{tabularx}
{\linewidth}{>{\hspace{-\tabcolsep}}l<{\hspace{-\tabcolsep}}
>{\hspace{-\tabcolsep}\enspace}llll
>{\raggedright\arraybackslash}X}
\lsptoprule
%\tbx{10-80}	&	*ə	&	PMP/PAMS	&	Hoti	&	gloss	&		\\	\midrule
%\endfirsthead
\tbx{10-80}	&	*ə	&	PMP/PAMS	&	Hoti	&	gloss	&		\\	\midrule
%\endhead
	&	{\hr}a	&	*tiR\tbr{ə}m	&	tih\tbr{a}l-a	&	oyster	&	\hp{\,}	\\
	&	{\hr}a	&	*mant\tbr{ə}R	&	mal\tbr{a}	&	cuscus	&	\hp{\,}	\\
	&	{\hr}a	&	*liq\tbr{ə}R	&	ri\tbr{a}	&	voice	&	\hp{\,}	\\
	&	{\hr}a	&	**vəl\tbr{ə}'	&	ɸel\tbr{a}ʔ-a	&	oil	&	\hp{\,}	\\
	&	{\hr}a	&	**təbəl\tbr{ə}	&	tibal\tbr{a}	&	snake	&	\hp{\,}	\\
	&	{\hr}a	&	**ma-səl\tbr{ə}	&	misor\tbr{a}	&	sick	&	\hp{\,}	\\
	&	{\hr}o	&	*bəqb\tbr{ə}q	&	ɸoɸ\tbr{o}-n	&	mouth	&	\citep{Ludeking1868}	\\
	&	{\hr}o	&	*kəd\tbr{ə}ŋ	&	ol\tbr{o}	&	stand	&	Bobot \it{a-kolan}, Masiwang \it{kolan}	\\
	&	{\hr}o	&	*kət\tbr{ə}p	&	ok\tbr{o}	&	bite	&	\hp{\,}	\\
	&	{\hr}o	&	*təkt\tbr{ə}k	&	tot\tbr{o}	&	cut off	&	\hp{\,}	\\
	&	{\hr}i	&	**bər\tbr{ə}n	&	boh\tbr{i}l-a	&	Nipa palm	&	\hp{\,}	\\
	&	{\hr}i	&	*dal\tbr{ə}m	&	lar\tbr{i}-	&	inside	&	Bobot \it{lala}, Masiwang \it{lalan}	\\
	&	{\hr}i	&	*las\tbr{ə}R	&	lis\tbr{i}-n	&	testicles	&	\hp{\,}	\\
	&	{\hr}e	&	*pus\tbr{ə}j	&	us\tbr{e}-	&	navel	&	\hp{\,}	\\
	&	{\hr}{\0}	&	*wahiy\tbr{ə}R	&	wal-a	&	water	&	\hp{\,}	\\
	&	{\hr}{\0}	&	**bat\tbr{ə}	&	bat-a	&	mango	&	\hp{\,}	\\
	&	{\hr}{\0}	&	**tabəl\tbr{ə}	&	tibon-a	&	bamboo	&	\hp{\,}	\\
\lspbottomrule
\end{tabularx}\mbox{}\\}

One change that is unique to Hoti in this eastern region is the
split of *t > \it{t, k}.
No conditioning factors have been identified,
apart from the fact that *t = \it{t} is the usual reflex in word-initial position.
Examples are given in \qf{10-81}.

\noindent{\wg\st{6pt}\tblr{>{\hspace{-\tabcolsep}}l<{\hspace{-\tabcolsep}}
>{\hspace{-\tabcolsep}\enspace}
lllll}
\lsptoprule
%\tbx{10-81}	&	*t	&	PMP/PAMS	&	Hoti	&	gloss	\\	\midrule
%\endfirsthead
\tbx{10-81}	&	*t	&	PMP/PAMS	&	Hoti	&	gloss	&\hp{\,}\\	\midrule
%\endhead
	&	{\hr}t	&	*\tbr{t}iRəm	&	\tbr{t}ihal-a	&	oyster	&\hp{\,}\\
	&	{\hr}t	&	*\tbr{t}əktək	&	\tbr{t}oto	&	cut off&\hp{\,}	\\
	&	{\hr}t	&	*\tbr{t}əbuh	&	\tbr{t}oɸ-a	&	sugarcane	&\hp{\,}\\
	&	{\hr}t	&	*\tbr{t}aliŋa	&	\tbr{t}ira	&	ear	&\hp{\,}\\
	&	{\hr}t	&	*\tbr{t}asik	&	\tbr{t}asiʔ-a	&	salt	&\hp{\,}\\
	&	{\hr}t	&	*\tbr{t}uak	&	\tbr{t}uaʔ	&	palm wine	&\hp{\,}\\
	&	{\hr}t	&	*ma\tbr{t}ay	&	ma\tbr{t}i-a	&	dead & (unexpected *-ay > \it{i})	\\
	&	{\hr}t	&	*ma-pu\tbr{t}iq	&	bi\tbr{t}-a / biti	&	white	\\
	&	{\hr}t	&	*uRa\tbr{t}	&	uʔa\tbr{t}-a / uat-a	&	vein, cord	&\hp{\,}\\
	&	{\hr}k	&	*\tbr{t}aqi	&	\tbr{k}ai-a	&	faeces	&\hp{\,}\\
	&	{\hr}k	&	*kə\tbr{t}əp	&	o\tbr{k}o	&	bite	&\hp{\,}\\
	&	{\hr}k	&	*qə\tbr{t}ut	&	o\tbr{k}it	&	fart	&\hp{\,}\\
	&	{\hr}k	&	*ba\tbr{t}aŋ	&	ɸa\tbr{k}a-	&	trunk	&\hp{\,}\\
	&	{\hr}k	&	*bi\tbr{t}uka	&	ɸa\tbr{k}ua	&	intestines	&\hp{\,}\\
	&	{\hr}k	&	*ma\tbr{t}a	&	ma\tbr{k}a-	&	eye &\citep{Ludeking1868}	\\
\lspbottomrule
\end{tabular}\mbox{}\\}